% haptics_feedback_node_explanation_rev.tex
\documentclass[a4paper,11pt]{article}
\usepackage{amsmath,amssymb}
\usepackage{bm}
\usepackage{siunitx}
\usepackage{hyperref}
\usepackage{enumitem}
\usepackage{physics}
\usepackage{booktabs}
\usepackage{geometry}
\geometry{margin=25mm}

\title{オペレータモデルを含むエネルギー保存則によるハプティクスフィードバック生成法 \\（修正版）}
\author{（原稿に基づき最小修正を適用）}
\date{}

\begin{document}
\maketitle

\section*{注意（符号と定義の統一）}
本稿では式の混乱を避けるために以下の符号規約を明確にする。

\begin{itemize}
  \item $\bm{w}\in\mathbb{R}^6$ は「\emph{そのポートが\textbf{システムに対して作用するレンチ}（すなわちシステムが受け取る力）}」を表す。
  \item $\bm{v}\in\mathbb{R}^6$ はそのポートのツイスト（並進速度・角速度）を表す（並進は m, 並進速度は m/s、回転は rad, 角速度は rad/s）。
  \item ポートに\emph{流入}（システムが受け取る）するパワーは
  \[
    P_{\text{port}} \;=\; \bm{w}^\top \bm{v}.
  \]
  \item 以降、$P_{\text{env}}$ や $P_{\text{op}}$ などは常に「システムへ流入するパワー」を意味することとする。
\end{itemize}

上の定義により以降の式は一貫して記述される（元稿で混在していた $P_{in}/P_{out}$ 表記は廃止し、すべて「システムへ流入するパワー」で統一した）。

\section{エネルギー収支（1ステップ予測）}
ロボット（外界ポート）とデバイス（オペレータポート）を含むシステムの 1 サイクル（ステップ $i\to i+1$）での外部からの流入／流出の予測は、次のように書く。

\begin{align}
  P_{\text{env}}(i) &= \bm{w}_{\text{measured}}(i)^\top \bm{v}_{\text{robot}}(i), \\
  P_{\text{op}}(i)  &= \bm{w}_{\text{op,est}}(i)^\top \bm{v}_{\text{op}}(i),
\end{align}
ここでいずれも「システムへ流入するパワー」であることに注意する。\\

予測されるシステムエネルギー変化（1ステップ分、$\Delta t$）は
\[
  \Delta E_{\text{pred}}(i) \;=\; \big( P_{\text{env}}(i) + P_{\text{op}}(i) - P_{\text{out,sys}}(i) \big)\Delta t,
\]
と書けるが、実装上は「オペレータへ提示するレンチにより予測される流出分」と「オペレータが逆に行う仕事」に着目して簡潔に扱う。以降は以下の簡便形で統一する：

\[
  \Delta E_{\text{pred}}(i) = \big( \bm{w}_{\text{nom}}(i)^\top \bm{v}_{\text{hand,pred}}(i+1) - \bm{w}_{\text{op,est}}(i)^\top \bm{v}_{\text{op}}(i) \big)\Delta t.
\]

\section{エネルギータンク更新（明確化）}
システムが内部に保持する利用可能エネルギーを $E_{\text{tank}}(i)$，その上限を $E_{\max}^{\text{tank}}$ とする。更新規則は符号をわかりやすく次の形で示す（常に「新しいタンク量 = 旧タンク量 + 系への流入量（正） - 系からの流出量（正）」の観点）：

\[
  E_{\text{tank}}(i+1) \;=\; \mathrm{clip}\big( E_{\text{tank}}(i) + \Delta E_{\text{toTank}}(i),\; 0,\; E_{\max}^{\text{tank}} \big),
\]
ここで $\Delta E_{\text{toTank}}(i)$ は「そのステップでタンクに蓄えられる（正）／取り出される（負）エネルギー」を意味する。実装上、\(\Delta E_{\text{toTank}}(i)= -\Delta E_{\text{pred}}(i)\) と扱う（※符号は上の定義に従って解釈すること）。\\

特別な場合の取り扱い：
\begin{itemize}
  \item $\Delta E_{\text{pred}}(i) > 0$（予測では流出）：まず $0 < \Delta E_{\text{pred}} \le E_{\text{tank}}$ なら通常提示（$ \bm{w}_{\text{feedback}} = \bm{w}_{\text{nom}}$）でタンクを取り崩す。
  \item $\Delta E_{\text{pred}}(i) > E_{\text{tank}}$（タンクより大きく流出する予測）：スケーリング係数 $\alpha\in[0,1]$ を用いて提示レンチを縮小し、流出をタンク以下に抑える（下節参照）。
  \item $\Delta E_{\text{pred}}(i) < 0$（予測では流入）：まずタンクに蓄えるが，蓄えきれない場合はダンパ追加で余剰を消費する（下節参照）。
\end{itemize}

\section{スケーリング（$\alpha$）による流出抑制 --- 明確化と数値安全化}
$\Delta E_{\text{pred}}(i) > E_{\text{tank}}(i)$ の時、提示レンチをスケールして流出を抑える。正しい代数式は次の通り：

\[
  \Delta E_{\text{pred}}(i) = \big( \alpha \,\bm{w}_{\text{nom}}(i)^\top \bm{v}_{\text{hand,pred}}(i+1) - \bm{w}_{\text{op,est}}(i)^\top \bm{v}_{\text{op}}(i) \big)\Delta t.
\]

これを $\alpha$ について解くと

\[
  \alpha = \frac{\dfrac{E_{\text{tank}}(i)}{\Delta t} + \bm{w}_{\text{op,est}}(i)^\top \bm{v}_{\text{op}}(i)}
               {\bm{w}_{\text{nom}}(i)^\top \bm{v}_{\text{hand,pred}}(i+1)}.
\]

\paragraph{数値的注意点（必須）}
\begin{enumerate}[label=(\alph*)]
  \item 分母が小さい／ゼロに近い場合は直接計算しない。判定閾値 $\delta_{\text{denom}}>0$（設計変数）を用い、\\
    \(\;|\bm{w}_{\text{nom}}^\top \bm{v}_{\text{hand,pred}}| < \delta_{\text{denom}}\Rightarrow \alpha \leftarrow 0\)（またはダンパへフォールバック）。
  \item 計算後は常に $\alpha \leftarrow \mathrm{clip}(\alpha,0,1)$。
  \item $\alpha$ の急激変化を避けるために1次遅れフィルタ（時定数 $\tau_\alpha$）を適用することを推奨する。
\end{enumerate}

\section{ダンパ追加によるエネルギー散逸（$D_{\text{add}}$）}
タンクに収めきれない流入がある場合、追加ダンパ $D_{\text{add}}=\mathrm{diag}(d_j)>0$ を導入して消費する。必要消費電力 $P_{\text{req}}$ を以下で定義する（単位：W）：

\[
  P_{\text{req}} = \frac{ \max\big(0,\; |\Delta E_{\text{nom}}| - \big(E_{\max}^{\text{tank}} - E_{\text{tank}}(i)\big)\big)}{\Delta t},
\]
ここで $\Delta E_{\text{nom}}$ はダンパ無しで予測される流入（負値の絶対）である。\\

各軸へ割当てるパワーの重み（正規化）を次で与える（数値安定化を含む）：
\[
  p_j = \max\big(0,\; w_{\text{nom},j}\, v_j\big),\qquad
  \hat q_j = \frac{p_j}{\sum_k p_k + \delta_q},
\]
ここで $\delta_q>0$ は分母ゼロ回避の閾値。\\

このとき各軸のダンピング係数は
\[
  d_j = \frac{P_{\text{req}}\,\hat q_j}{v_j^2 + \varepsilon_j},
\]
ただし実装上は必ず上限クリップ $d_j \leftarrow \min(d_j,\, d_{j,\max})$ を行う。\\

\paragraph{設計上の注意}
\begin{itemize}
  \item $\varepsilon_j$ は速度スケールに合わせて選ぶ（例：速度スケール $v_{\text{scale}}\!=\!1\times10^{-2}$ の場合 $\varepsilon \approx v_{\text{scale}}^2 = 1\times10^{-4}$）。極端に小さい値（$10^{-8}$ 等）は実装では不安定になる場合があるため注意。
  \item $d_{j,\max}$ はデバイスの最大力 $F_{j,\max}$ と最小速度想定 $v_{j,\min}$ より $d_{j,\max}\approx F_{j,\max}/v_{j,\min}$ を基に決め、安全余裕を取ること。
\end{itemize}

\section{オペレータインピーダンス同定（RLS）の明示}
逐次最小二乗（RLS）でオペレータの等価インピーダンス $\theta_i=[M_i,B_i,K_i]^\top$ を推定する際、数式と観測値の定義を明確化する。

\paragraph{回帰式}
各成分について
\[
  F_i \approx \phi_i^\top \theta_i,\qquad
  \phi_i = \begin{bmatrix} \ddot x_i \\ \dot x_i \\ x_i - x_{\text{ref}} \end{bmatrix},
\]
ただし $x$ は成分ごと（並進3, 回転3）で扱う（回転は rad 単位で表す）。

\paragraph{RLS 更新（忘却因子あり）}
\begin{align}
  Q_i &= \frac{1}{\lambda}\left( Q_{i-1} - \frac{Q_{i-1}\phi_i\phi_i^\top Q_{i-1}}{\lambda + \phi_i^\top Q_{i-1}\phi_i} \right), \\
  \theta_i &= \theta_{i-1} + Q_i \phi_i \big(y_i - \phi_i^\top \theta_{i-1}\big),
\end{align}
ここで $\lambda\in(0,1]$ は忘却係数（設計変数）であり、$Q_0$ は初期共分散（行列）である。\\

\paragraph{観測値 $y_i$ の定義（重要）}
元稿では $y_i$ を「前ステップに提示したレンチ」としているが、\textbf{実装上はデバイスで\emph{実際に発生した}力/トルク（実測）を使うべき}である。コマンド値をそのまま $y_i$ に用いるとアクチュエータ飽和や非線形の影響を同定値に取り込んでしまう。従って $y_i$ は「デバイスのセンサで計測した反力（または実出力の推定）」を用いることを明記する。

\section{6次元（並進/回転）の単位と分離}
6自由度を要素ごとに扱う場合、\emph{並進成分}と\emph{回転成分}は物理単位が異なるため次の点を必ず明記する：
\begin{itemize}
  \item 並進の質量 $M_{\text{trans}}$ は kg、回転の慣性 $M_{\text{rot}}$ は kg$\cdot$m$^2$ とする。
  \item ダンピング $B_{\text{trans}}$ は Ns/m、回転は Nms/rad。
  \item 剛性 $K_{\text{trans}}$ は N/m、回転は Nm/rad。
  \item それぞれ別個に初期値レンジを与えること（下節で推奨オーダーを示す）。
\end{itemize}

\section{設計変数一覧（追補）と推奨オーダー}
以下は原稿に含まれる変数に加えて、実装前に設計者が\textbf{必ず}決めておくべき変数と初期推奨オーダー（目安）である。数値は典型的ハプティクス実験を想定した目安であり、デバイス特性により調整が必要である。

\begin{table}[h]
\centering
\begin{tabular}{l l l}
\toprule
変数 & 単位 & 推奨オーダー／備考 \\
\midrule
$\Delta t$ & s & 0.001--0.01（ハプティクスなら1~\si{ms}推奨） \\
$E_{\max}^{\text{tank}}$ & J & 1--50（デスクトップ:1--10 J, 産業:10--50 J） \\
$E_{\text{tank}}(0)$ & J & 0--$E_{\max}$（安全優先なら小さめ） \\
$k_{\text{scale,trans}}$ & - & 0.2--0.5（並進） \\
$k_{\text{scale,rot}}$ & - & 0.05--0.2（回転） \\
$M_{0,\text{trans}}$ & kg & 0.1--5 (典型 0.5--2) \\
$M_{0,\text{rot}}$ & kg$\cdot$m$^2$ & $10^{-4}$ -- $10^{-1}$ (典型 $10^{-3}$) \\
$B_{0,\text{trans}}$ & Ns/m & 0.1--50 (典型 1--20) \\
$B_{0,\text{rot}}$ & Nms/rad & $10^{-3}$ -- 1 (典型 0.01--0.5) \\
$K_{0,\text{trans}}$ & N/m & 10--1000 (典型 50--500) \\
$K_{0,\text{rot}}$ & Nm/rad & 0.1--100 (典型 1--20) \\
$\lambda$ (RLS) & - & 0.98--0.9999（$\Delta t$ に依存。1~\si{ms}で0.999なら有効記憶約1s） \\
$Q_0$ & - & 対角（大きめ）：例 diag(100,100,100)（スケール依存） \\
$\varepsilon_j$ & (速度)$^2$ & $v_{\text{scale}}^2$（例 $v_{\text{scale}}=10^{-2}$ $\Rightarrow \varepsilon=10^{-4}$） \\
$\delta_{\text{denom}}$ & W & small positive（例 $10^{-6}$--$10^{-3}$） \\
$\delta_q$ & W & small positive（分母回避。例 $10^{-6}$--$10^{-3}$） \\
$d_{j,\max}$ & Ns/m または Nms/rad & デバイスの $F_{j,\max}/v_{j,\min}$ を基に決定、さらに安全余裕を取る \\
$F_{j,\max},\;T_{j,\max}$ & N, Nm & デバイス性能に依存（必須） \\
\bottomrule
\end{tabular}
\end{table}

\section{数値安全化・実装チェックリスト}
実装前に必ず確認する項目を列挙する。
\begin{enumerate}
  \item 符号規約（本文冒頭）に基づくすべての式が一貫しているか。
  \item $\alpha$ 計算の分母小判定（$\delta_{\text{denom}}$）とクリッピングが入っているか。
  \item $d_j$ の $\varepsilon_j$ と $d_{j,\max}$ が適切に設定されているか。
  \item RLS の $y_i$ を「実測力（デバイス）」に変更しているか。$Q_0,\lambda$ が定義されているか。
  \item 並進/回転で単位が分離され、初期パラメータがそれぞれ与えられているか。
  \item センサ／コマンドの遅延を考慮した時間整合（タイムスタンプ整列）が入っているか。
  \item 異常値（センサ外れ値、 saturations）でのフォールバック処理を用意しているか。
\end{enumerate}

\section*{補遺：最小変更の意図}
本修正版は元稿の構成・記述を保ちつつ，以下の最小限の修正を加えました：
\begin{itemize}
  \item 「パワー（エネルギー）の符号規約」を文頭に明示して式の混在を解消。
  \item $\alpha$ の代数式を明確に括弧で示し，分母ゼロ回避とクリッピング／ローパスを付加。
  \item RLS の更新式を「忘却因子ありの標準形」に統一し，観測 $y_i$ の定義を「実測力」に修正。
  \item ダンパ割当てでの $\varepsilon_j$, $\delta_q$, $d_{j,\max}$ 等の数値安全化を明記。
  \item 6次元（並進/回転）の単位分離と具体的な初期レンジの目安を追加。
\end{itemize}

\vspace{1em}
\noindent
（この文書は元の TeX を基に最小修正を行ったものである。参照元原稿：:contentReference[oaicite:1]{index=1}）

\end{document}
